%% =====================================================================
%%  B1-T-05-02 -- what B1 contributes to "Follow the Interest"
%%  -------------------------------------------------------------------
%%  LAYER A: APPEND-ONLY. Written once, then left alone. Material found
%%  only here sits in the `unique` slot and is protected by rule R10.
%% =====================================================================
%% @kind: source
%% @id: B1-T-05-02
%% @book: B1
%% @topic: T-05-02
%% @locator: Tactic 8, pp. 89--108; Tactic 3, pp. 31--42
%% @status: distilled
%% @unique: yes
%% @integrated: yes
%% @updated: 2026-09-19

\dossier{B1}{T-05-02}{\bookshort{B1}}{Tactic 8, pp. 89--108; Tactic 3, pp. 31--42}

\begin{distilled}
  \item Avoiding being taken in financial dealings is reduced to a single instruction: watch where the money goes.
  \item A person proposing a deal must have something to gain from it, or they would not spend the breath; the more they emphasise your gain, the more they stand to make at your expense.
  \item Favourable terms are said to speak for themselves and to be offered with a take-it-or-leave-it stance, not argued at length.
  \item Favouritism is treated as a standing feature of the world rather than an aberration, and therefore as something to understand and use rather than to complain about.
\end{distilled}

\begin{unique}
  \item The volume rule: the amount of emphasis placed on your benefit is inversely related to how favourable the terms actually are.
  \item Favouritism as a normal structural feature of institutions rather than a corruption of them.
\end{unique}
